\begin{lemma}
	For any graph with no negative weight cycle, $\exists p$, s.t. $\ell_p(u,v) \ge 0$ for any $(u,v) \in E$.
\end{lemma}
\begin{proof}
	We add a new vertex $s$ and connect it to all other vertices with zero weight edge. Then we run Bellman-Ford algorithm on this new graph $G'$ and let $p(v) = \text{dist}_{G'}(s,v)$.
	As $\dist(s,v) \le \dist(s,u) + \ell(u,v)$, we have $\ell_p(u,v) = \ell(u,v) + p(u) - p(v) = \ell(u,v) + \dist(s,u) - \dist(s,v) \ge 0$.
\end{proof}


\subsection{Dijkstra's Algorithm}

Now assume $\ell: E \to \mathbb{R}^+ \cup \{0\}$, we can use Dijkstra's algorithm to find the shortest path from a source vertex to all other vertices in $O(m + n \log n)$ time.

Denote $\dist_S(s,v)$ as the distance from $s$ to $v$ in the subgraph induced by $S\cup \{v\}$, we have the following claim and observation:

\begin{lemma}
	If $v\not\in S$ minimizes $\dist_S(s,v)$, then $\dist_S(s,v) = \dist(s,v)$.
\end{lemma}
\begin{proof}
	Let $P$ be the shortest path from $s$ to $v$. If $P$ only contains vertices in $S\cup \{v\}$, then $\dist_S(s,v) = \dist(s,v)$ and we are done. Otherwise, let $u$ be the first vertex in $P$ that is not in $S$. Then the subpath from $s$ to $u$ only contains vertices in $S$, so $\dist_S(s,u) = \dist(s,u)$. As $\dist(u,v)\ge 0$ and $v$ minimizes $\dist_S(s,v)$, we have $\dist_S(s,v) \le \dist_S(s,u) + \dist(u,v) = \dist(s,u) + \dist(u,v) = \dist(s,v)$. On the other hand, as $S\cup \{v\}$ is a subgraph of $G$, we have $\dist_S(s,v) \ge \dist(s,v)$. Therefore, $\dist_S(s,v) = \dist(s,v)$.
\end{proof}

\begin{corollary}
	If we list vertices in increasing order of their distance from $s$, then for any vertex $v$, the shortest path from $s$ to $v$ only contains vertices that are listed before $v$.
\end{corollary}

Therefore, we can start with an empty $S$ and iteratively add the vertex that minimizes $\dist_S(s,v)$ until all vertices are added, with the loop invariant on $d(v)$ that we maintain:

\begin{itemize}
	\item If $v\in S$, then $d(v) = \dist(s,v)$.
	\item If $v\not\in S$, then $d(v) = \dist_S(s,v)$.
\end{itemize}

\begin{algorithm}[H]
	\caption{Dijkstra's Algorithm}\label{alg:dijkstra}
	Initialize $S = \emptyset$ and $d(v) = \infty$ for all $v\in V\setminus \{s\}$, $d(s) = 0$\;
	Insert all vertices into a set $Q$ with key $d(v)$\;
	\While {$Q$ is not empty}{
		$v \gets $ vertex in $Q$ with minimum key, $Q$ delete $v$\;
		$S \gets S\cup \{v\}$\;
		For each edge $(v,u)\in E \land u\in Q$, $d(u) \gets \min(d(u), d(v) + \ell(v,u))$\;
	}
	\Return{$d$}
\end{algorithm}

The set $Q$ should be a priority queue that supports the following operations:
\begin{itemize}
	\item insert$(x,k)$: $Q:=Q\cup \{x\}, d(x) := k$.
	\item deletemin: return $x$ with minimum $d(x)$ and $Q := Q\setminus \{x\}$.
	\item decreasekey$(x,k)$: $d(x) := \min\{d(x), k\}$.
\end{itemize}

If we maintain the priority queue with a binary heap, the time complexity is $O(m \log n)$. This is because insert, deletemin and decreasekey operations all take $O(\log n)$ time.
However, here we present a new data structure that has amortized $O(1)$ time for insert and decreasekey, and $O(\log n)$ time for deletemin, that leads to a time complexity of $O(m + n \log n)$ for Dijkstra's algorithm.

\begin{remark}[Solving APSP]

	We can't speed up general SSSP with negative weight edges by using Dijkstra's algorithm with Fibonacci heap, as calculating price function is essentially running Bellman-Ford algorithm, which takes $O(mn)$ time.

	However, for APSP, we can first calculate the price function $p$ in $O(mn)$ time, then reweight the edges and run Dijkstra's algorithm from each vertex, which takes $O(n(m + n \log n))$ time. Therefore, the total time complexity is $O(mn + n^2 \log n)$, which is better than $O(n^3)$.

\end{remark}

\section{Fibonacci Heap}

To start with, we need to define some concepts on rooted trees.

\begin{definition}[Heap-ordered Tree]
	A rooted tree $T$ is heap-ordered if for any node $u$ and its parent $v$, we have $d(v) \le d(u)$.
\end{definition}
\begin{note}
	Two heap-ordered trees $T_1$ and $T_2$ can be merged by making the root of the tree with larger key the child of the root of the tree with smaller key. This operation takes $O(1)$ time.
\end{note}

\begin{definition}[Binomial Tree]
	A binomial tree $B_k$ with rank $k$ is defined recursively as follows:
	\begin{itemize}
		\item $B_0$ is a single node.
		\item $B_k$ consists of two $B_{k-1}$ trees, where the root of one $B_{k-1}$ is the rightmost child of the root of the other $B_{k-1}$.
	\end{itemize}
\end{definition}

For example, $B_4$ is shown below:

\begin{center}
	\begin{forest}
		for tree={
		circle,             % Shape of the nodes
		draw,               % Draw the border
		minimum size=2em,   % Size of the nodes
		inner sep=1pt,
		math content,       % Allows using math mode inside nodes automatically
		s sep=5mm,          % Minimum horizontal distance between siblings
		l sep=10mm,         % Vertical distance between levels
		edge={semithick}    % Thickness of the connecting lines
		}
		% Root of B_4
		[B_4
		% Child: B_0
		[B_0]
		% Child: B_1
		[B_1
		[B_0]
		]
		% Child: B_2
		[B_2
		[B_0]
		[B_1 [B_0]]
		]
		% Child: B_3
		[B_3
		% Child of B_3: B_0
		[B_0]
		% Child of B_3: B_1
		[B_1 [B_0]]
		% Child of B_3: B_2
		[B_2
		[B_0]
		[B_1 [B_0]]
		]
		]
		]
	\end{forest}
\end{center}

One can find that there're $\binom{k}{i}$ nodes of depth $i$ in $B_k$, and the total number of nodes in $B_k$ is $\sum_{i=0}^k \binom{k}{i} = 2^k$ (that's why it's called a binomial tree.)
Moreover, the subtree rooted at the $i$-th child of the root of $B_k$ is a $B_{i-1}$ tree.

Now we collect multiple binomial trees into a binomial heap:

\begin{definition}[Binomial Heap]
	A binomial heap is a collection of heap-ordered binomial trees, where there is at most one binomial tree of any rank.
\end{definition}

In this data structure, insert and deletemin are easy to implement.

For insert$(x,k)$, we
\begin{itemize}
	\item Create a new $B_0$ tree that contains a single node $x$ with key $k$.
	\item Merge $H$ with the current binomial heap, if there are two trees of the same rank, we merge them into a tree of higher rank until there is at most one tree of any rank.
\end{itemize}

For deletemin, we
\begin{itemize}
	\item Find the tree $B_k$ with minimum root key, and remove the root $x$ from $B_k$.
	\item Let $H'$ be the binomial heap formed by the children of $x$, we merge $H'$ with the current binomial heap, and link the trees of the same rank until there is at most one tree of any rank.
\end{itemize}

\begin{lemma}
	insert and deletemin operations on a binomial heap with $n$ nodes take $O(\log n)$ time.
\end{lemma}
\begin{proof}
	All sub-operations are $O(1)$, so the operation is dominated by the number of linkings.
	As every linking reduces the number of trees by one, and there are at most $O(\log n)$ trees in a binomial heap with $n$ nodes, the number of linkings is at most $O(\log n)$.
\end{proof}

However this structure does not support decreasekey operation efficiently. To solve this problem, we introduce the Fibonacci heap, which is a more relaxed version of the binomial heap.

\begin{definition}[Fibonacci Heap]
	A Fibonacci heap is a collection of heap-ordered trees (not necessarily binomial), with at most one tree per rank.
\end{definition}

While the insert and deletemin operations are the same as in a binomial heap, the decreasekey operation is implemented as follows:

\begin{itemize}
	\item Let $x$ be the node whose key is decreased, and $y$ be the parent of $x$.
	\item Cut the edge between $x$ and $y$, and make $x$ a new tree in the heap.
	\item Mark $y$. If $y$ is already marked (i.e. it lost 2 children), we cut the edge between $y$ and its parent (and mark the parent as well), and make $y$ a new tree in the heap. We repeat this process until we reach a root or an unmarked node.
\end{itemize}

\begin{definition}
	The rank of a tree (or the root of the tree) is $k$ if its root has $k$ children, or it has $k-1$ children and is marked.
\end{definition}

\begin{lemma}
	The $i$-th child of a node with rank $k$ has rank at least $i-1$.
\end{lemma}
\begin{proof}
	We prove by induction on $i$. For $i=1$, the claim is trivial. Assume the claim holds for $i-1$, we want to show that the $i$-th child of a node with rank $k$ has rank at least $i-1$.

	When the $i$-th child is linked to the root, the root has at least $i-1$ children, so its rank is at least $i-1$. When a child is cut from its parent, the rank `requirement' of the other children will only be looser, so the $i$-th child will still have rank at least $i-1$.

	In all, the lemma holds.
\end{proof}

Denote $T_k$ to be the minimum size of a tree with rank $k$, we have $T_0 = 1, T_1 = 2$, and $T_k \ge 1 + T_0 + T_1 + \cdots + T_{k-3} + T_{k-2}$ for $k\ge 2$ (as we can't guarantee whether that last child exists), which leads to $T_k \ge T_{k-1} + T_{k-2}$ for $k\ge 2$.

One can find that $T_k$ is the Fibonacci sequence, so $T_k \ge F_k \ge \phi^{k+1}$, where $\phi = \frac{1+\sqrt{5}}{2}$ is the golden ratio. Therefore, we have:

\begin{corollary}
	The maximum rank of a tree in a Fibonacci heap with $n$ nodes is at most $O(\log n)$.
\end{corollary}

Now we are ready to analyze the time complexity of the operations on a Fibonacci heap.

\begin{theorem}
	The insert and decreasekey operations on a Fibonacci heap with $n$ nodes take $O(1)$ amortized time, and the deletemin operation takes $O(\log n)$ amortized time.
\end{theorem}
\begin{proof}
	We start by defining a potential function. Note that:
	\begin{itemize}
		\item Linking reduces the number of trees.
		\item Cutting increases the number of trees, but decreases the number of marked nodes (or increase with an upper bound of 1).
	\end{itemize}
	Therefore, we can let the potential function be $\Phi = t + 2m$, where $t$ is the number of trees and $m$ is the number of marked nodes in the heap. The amortized cost of any operation is then $a_i = t_i + \Delta_i$, with $\Delta_i$ being the change of the potential function.

	Suppose the $i$-th operation is insert$(x,k)$, we create a new tree with $x$ as the root, then link the trees that have the same rank, so $t_i = \text{\# linkings} + 1$. On the other hand, each linking reduces the number of trees by one, so $\Delta_i = 1 - \text{\# linkings}$ (the 1 comes from the new tree), therefore, $a_i = t_i + \Delta_i = 2$.

	Suppose the $i$-th operation is decreasekey$(x,k)$, we cut $x$ from its parent, then recursively cut the parents if they're marked, and link the trees with the same rank, so $t_i = \text{\# linkings} + \text{\# cuts}$. Let the number of cuts be $c$, then the number of trees increases by $c$ and then decreases by the number of linkings, and the number of marked nodes decreases by $c-2$ (increases if $c-2$ is negative.) This means $\Delta_i = \text{\# cuts} - \text{\# linkings} - 2(\text{\# cuts} - 2) = 4 - \text{\# linkings} - \text{\# cuts}$, therefore, $a_i = t_i + \Delta_i = 4$.

	Suppose the $i$-th operation is deletemin, we find the tree with minimum root key, remove the root $x$ from the tree, and link the trees with the same rank. Let $k$ be the rank of the tree, then we have $t_i = O(k) + \text{\# linkings}$ (the $O(k)$ comes finding the minimum), and $\Delta_i = O(k) - \text{\# linkings}$, therefore, $a_i = t_i + \Delta_i = O(k) = O(\log n)$.

	This completes the proof.
\end{proof}
































